% TOP_PROBLEM: {"id": 20002527, "title": "The L2-signature of toric hyperkahler quotients", "category_id": 16, "category": {"id": 16, "name": "physics", "display_name": "Mathematical Physics", "description": "Problems at the intersection of mathematics and physics.", "slug": "physics", "order_index": 16, "created_at": "2026-07-31T15:26:25.671Z"}, "status": "partially_solved", "statusQualification": "The assertions about concentration in middle degree and self-duality are historical source context with geometric assumptions, not newly verified universal statements. This is a third first step in the shared programme of 20002525–20002527. The catalogue title/status was paired with a machine-generated research attempt; that attempt is not peer-reviewed evidence of settlement. The statement covers the original source questions rather than only the title’s subcase."}
% FIRST_POSTED: 2026-10-11
% ENTRY_KIND: statement_only
% SUBMISSION: {"kind": "statement_only", "category": "top_problems", "problem_id": "20002527", "proposal_type": "established", "source": "American Institute of Mathematics, $L^2$ harmonic forms in geometry and string theory; notes by Anda Degeratu and Mark Haskins, March 2004, 21 March 2004, Hausel wrap-up programme and first step 3. https://aimath.org/WWN/l2harmonic/l2harmonic.pdf", "status_evidence": "Inherited source-catalogue status partially_solved preserved. Background has machine-generated research claims explicitly marked not peer reviewed; these are not independent status evidence. Primary source accessed 11 October 2026; compared against the supplied catalogue excerpt. Historical wording and partial-result remarks were checked for scope; no comprehensive current-literature or proof-verification audit was performed.", "status_checked_on": "2026-10-11", "duplicate_check": "Full 15,773-record source registry checked for identical normalized statements; available site registry checked for identical normalized title/statement and canonical-ID collision. Canonical ID is absent from the site registry. Related canonical records 20002525, 20002526. No mathematical-equivalence literature audit performed.", "attribution": "Tamás Hausel. Statement preparation: OpenAI GPT-6 Codex assistant, generation 11 October 2026; source-matching and notation editing by the assistant, no independent human mathematical verification."}
% SOURCE_STATEMENT: "3. (Hausel) We know that all the L2 Hodge cohomology of these spaces must be in the middle dimension and that any L2-harmonic form in middle dimension is ± self-dual. So it is enough to understand the\n\nL2-signature. This will have consequences for the sign of σL2 (X) and is related to combinatorics and matroids.\n\nMore general problem: Prove a suitable index theorem on such manifolds. Will there be contributions from subsytems??\n\nQuestions (G. Etesi):"
% Preparation and final self-audit: CONTRIBUTING.md prompts C and D; source/audit artifacts retained in the local physics_statements_2026_10_11 batch.
% Imported from: physics_statements.tex; UnsolvedMath ID 20002527
% !TeX program = lualatex
\documentclass[11pt]{article}
\usepackage[margin=25mm]{geometry}
\usepackage{fontspec}
\setmainfont{Latin Modern Roman}
\usepackage{amsmath,amssymb,mathtools}
\usepackage{unicode-math}
\setmathfont{Latin Modern Math}
\usepackage{xurl,hyperref}
\hypersetup{colorlinks=true,urlcolor=blue}
\setcounter{secnumdepth}{0}
\setlength{\parindent}{0pt}
\setlength{\parskip}{5pt}
\setlength{\emergencystretch}{3em}
\newcommand{\sref}[2]{\hyperlink{source:#1:#2}{[S#2]}}
\newcommand{\cataloguescope}[1]{\paragraph{Recorded target.}#1}
\newcommand{\R}{\mathbb R}
\newcommand{\C}{\mathbb C}
\newcommand{\N}{\mathbb N}
\newcommand{\Z}{\mathbb Z}
\newcommand{\Q}{\mathbb Q}
\newcommand{\F}{\mathbb F}
\newcommand{\eps}{\varepsilon}
\begin{document}
% UnsolvedMath ID 20002527; source catalogue code AIM-PHYSICS-0035
\setcounter{section}{20002526}
\section{The L2-signature of toric hyperkahler quotients}\label{20002527}
{\small UnsolvedMath ID 20002527 · Mathematical Physics}
\subsection{Definitions and mathematical statement}

\paragraph{Shared notation.}$\mathbb N=\{1,2,\ldots\}$; the standard number fields have their usual meanings. All additional conventions are stated locally.


\paragraph{Definitions and notation.}A hyper-Kähler metric has three parallel complex structures satisfying the quaternionic relations. In four dimensions a gravitational instanton is a complete noncompact Ricci-flat geometry in the source’s self-dual setting. ALE and ALF describe different asymptotic geometries; ALF has an asymptotic circle direction over a three-dimensional base. For a Riemannian manifold $(M,g)$, $L^2$ norms use the metric volume and pointwise metric on forms. Reduced $L^2$ cohomology is $\ker d/\overline{\operatorname{im}d}$ for the specified closed $L^2$ complex. On a complete manifold its harmonic representatives satisfy $d\alpha=d^*\alpha=0$ and $\int_M|\alpha|^2<\infty$. In even middle degree, the Hodge star splits harmonic forms into self-dual and anti-self-dual sectors. Their dimension difference is an $L^2$ signature when finite and suitably defined. An index theorem must specify the Fredholm operator and the asymptotic or boundary corrections needed on a noncompact space. Toric hyper-Kähler quotients have combinatorial data associated with the reducing torus.

\paragraph{Catalogue wording retained for traceability.}
For the source’s toric hyper-Kähler quotients, understand the $L^2$ signature and its sign, using the source’s middle-degree harmonic-form context and its relation to combinatorics and matroids. More generally, prove an appropriate index theorem for these manifolds and determine whether contributions from subsystems occur.
\paragraph{Scope and source qualification.}The assertions about concentration in middle degree and self-duality are historical source context with geometric assumptions, not newly verified universal statements. This is a third first step in the shared programme of 20002525–20002527. The catalogue title/status was paired with a machine-generated research attempt; that attempt is not peer-reviewed evidence of settlement. The statement covers the original source questions rather than only the title’s subcase.
\subsection{Short English statement}
Compute L² signatures and suitable index formulae for toric hyper-Kähler quotients, including possible subsystem terms.
\subsection{Sources}
\begin{itemize}
\item[\hypertarget{source:20002527:1}{S1}] ULAM.AI and the UnsolvedMath Contributors, UnsolvedMath, record 20002527 (AIM-PHYSICS-0035). Catalogue curation under CC BY 4.0; definitions and English exposition added locally.
\url{https://huggingface.co/datasets/ulamai/UnsolvedMath}.
\item[\hypertarget{source:20002527:2}{S2}] American Institute of Mathematics, $L^2$ harmonic forms in geometry and string theory; notes by Anda Degeratu and Mark Haskins, March 2004, 21 March 2004, Hausel wrap-up programme and first step 3.
\url{https://aimath.org/WWN/l2harmonic/l2harmonic.pdf}.
\end{itemize}
\label{end:20002527}
\end{document}
